\documentclass{article}

% TMLR format
\usepackage[preprint]{tmlr}

\newcommand{\openreview}{\url{https://openreview.net/forum?id=XXXX}}
\def\month{06}
\def\year{2026}

\usepackage{amsmath,amssymb,amsthm}
\usepackage{booktabs}
\usepackage{graphicx}
\usepackage{hyperref}
\usepackage{xcolor}
\usepackage{array}
\usepackage{enumitem}
\usepackage{multirow}

\newcommand{\Gr}{\mathrm{Gr}}
\newcommand{\IIA}{\mathrm{IIA}}
\newcommand{\R}{\mathbb{R}}
\newcommand{\Z}{\mathbb{Z}}
\newcommand{\dgr}{d_{\Gr}}
\newcommand{\Span}{\mathrm{span}}
\newcommand{\KL}{D_\mathrm{KL}}
\newcommand{\JS}{D_\mathrm{JS}}
\DeclareMathOperator{\diag}{diag}

\newtheorem{assumption}{Assumption}
\newtheorem{proposition}{Proposition}

\title{When Does Linear Causal Abstraction Work?\\
Mapping the Boundary on the Grassmannian}

\author{
Elliot Tower \\
\texttt{elliot@elliottower.com}
}

\begin{document}
\maketitle

\begin{abstract}
Distributed Alignment Search (DAS) identifies low-dimensional linear subspaces
that mediate model behavior under interchange interventions, implicitly assuming
that the relevant causal variable lives on the Grassmannian manifold
$\Gr(k,d)$. When does this assumption hold, and what happens when it fails?

We construct an atlas of 14 modular arithmetic operations and report four
findings. \textbf{(1)}~Operations partition into three sharp classes---always,
stochastic, and never Grassmannian---governed by grokking: Grassmannian
variables appear if and only if the model generalizes.
\textbf{(2)}~Linear DAS returns $\IIA = 0.0$ on fully grokked modular
addition, while memorized models achieve $\IIA \geq 0.86$ without geometric
structure, demonstrating that IIA alone cannot distinguish genuine causal
variables from lookup tables.
\textbf{(3)}~A Structured pi-SAE---a label-conditional sparse autoencoder
with end-to-end intervention training---recovers nonlinear causal variables
where DAS fails, achieving IIA~=~0.97 on GPT-2 hypernymy at $k{=}1$ (DAS:
0.07). Unconstrained nonlinear baselines achieve perfect IIA vacuously via
degenerate encoder-decoders; we introduce \emph{intervention faithfulness}
metrics that expose this failure mode.
\textbf{(4)}~Mechanistic analysis reveals that nonlinearity is necessary for
grokking but attention is not (DMF + ReLU groks 18$\times$ faster than the
full transformer); that Fourier structure emerges \emph{after} generalization,
not before it; that weight-space SVD predicts causal geometry without
optimization ($\IIA = 0.999$ vs.\ DAS's $0.150$); and that the causal
subspace identity is gauge freedom---10 seeds produce orthogonal subspaces
that all support equally valid interventions.
\end{abstract}

\section{Introduction}
\label{sec:intro}

Mechanistic interpretability seeks to identify causally active,
human-understandable variables inside neural networks. Causal abstraction
formalizes this goal: a model variable is ``causal'' if interventions on its
internal representation behave like interventions on a high-level algorithmic
variable \citep{geiger2021causal}. Distributed Alignment Search (DAS)
operationalizes causal abstraction by learning a rotation $Q \in \R^{d \times k}$
that identifies a $k$-dimensional subspace of the residual stream. Swapping the
in-subspace component between two inputs and measuring whether the model
produces the corresponding counterfactual output yields Interchange Intervention
Accuracy (IIA) \citep{geiger2023finding}.

DAS has been applied to syntax \citep{geiger2023finding}, arithmetic
\citep{wu2024interpretability}, and board games \citep{li2023emergent}, but
every application embeds a structural assumption: the causal variable is a
\emph{linear} subspace---a point on the Grassmannian manifold $\Gr(k,d)$. When
the true causal variable is linear, DAS recovers it. When it is not, DAS may
still return a high-IIA subspace that reflects memorization or distributional
artifacts rather than genuine causal structure.

We do not ask whether DAS ``works'' in the narrow sense of achieving high IIA---it
almost always does. We ask: \textbf{when does the linear assumption hold?} Is the
subspace DAS recovers a genuine Grassmannian point with the geometric properties
expected of a linear causal variable, or an artifact of searching in the wrong
function class? And when it fails, \textbf{can a nonlinear method recover the
causal variable?}

To answer these questions, we construct an atlas of 14 operations over $\Z_p$
spanning the boundary between linear and nonlinear causal structure. For each, we
train a single-layer transformer in the standard grokking setup, fit DAS at
multiple subspace dimensions, and evaluate the recovered subspace with three
diagnostics beyond IIA: equivariance under the operation's symmetry group, circle
geometry of the DAS projection, and intrinsic dimension from $k$-sweeps. We then
apply a Structured pi-SAE \citep{narayanaswamy2017learning} to the
never-Grassmannian operations, showing that nonlinear causal variables exist even
where linear methods fail.

The atlas reveals a sharp, grokking-governed partition into three classes:
seven operations \emph{always} produce Grassmannian variables, two are
\emph{stochastic} (same operation, different seeds, opposite outcomes), and
five \emph{never} do. Grokked modular addition returns $\IIA = 0.0$ under
linear DAS despite having learned the correct algorithm, while memorized
models achieve $\IIA \geq 0.86$ without geometric structure---demonstrating
that IIA alone cannot distinguish genuine causal variables from lookup tables.
A Structured pi-SAE recovers nonlinear causal variables where DAS fails,
achieving IIA~=~0.97 on GPT-2 hypernymy at $k{=}1$ (DAS: 0.07).

Beyond the atlas, we ask \textbf{how does grokking create causal structure?}
Architecture ablations show that nonlinearity is necessary but attention is
not. Fine-grained trajectory analysis reveals that Fourier structure emerges
\emph{after} generalization, and that weight-space SVD predicts causal
geometry without optimization. The causal subspace itself is gauge
freedom---the model does not prefer a canonical subspace.

\paragraph{Contributions.}
\begin{enumerate}[leftmargin=*]
\item \textbf{Atlas and stochastic grokking.} We construct an atlas of 14
  operations mapping the boundary of linear causal abstraction. We discover
  stochastic grokking---seed-dependent outcomes for operations near the
  boundary---providing the cleanest evidence that generalization, not
  architecture, enables Grassmannian causal variables.
\item \textbf{IIA insufficiency and faithfulness metrics.} We show that
  high IIA is compatible with memorization ($\IIA = 1.0$ at $k{=}4$ for
  non-grokking operations) and with degenerate encoder-decoders (NL-DAS
  achieves $\IIA = 1.0$ via lookup tables). We introduce equivariance
  testing, intervention diversity ratio, and continuous distributional
  measures (KL, JS) as necessary complements.
\item \textbf{Structured pi-SAE.} A label-conditional sparse autoencoder
  with end-to-end intervention training recovers nonlinear causal
  variables where DAS fails entirely, with iVAE identifiability guarantees.
  Unconstrained nonlinear baselines are vacuous; the ELBO prevents
  degeneracy.
\item \textbf{Mechanistic analysis.} Nonlinearity is necessary for grokking
  but attention is not (DMF + ReLU groks 18$\times$ faster). Fourier
  structure lags generalization by $\sim$1{,}000 epochs. Weight-space SVD
  achieves $\IIA = 0.999$ without optimization. The causal subspace identity
  is gauge freedom (10 seeds, pairwise overlap $\approx 0.008$). Spectral
  profiling works on toy models but fails under superposition.
\item \textbf{Evaluation hierarchy.} We propose a four-level framework---interchange effectiveness, interchange faithfulness, distributional
  fidelity, and structural diagnostics---for evaluating causal variable
  claims.
\end{enumerate}


\section{Background}
\label{sec:background}

\subsection{Causal abstraction and DAS}

Let $h(x) \in \R^d$ be an intermediate activation for input $x$. DAS searches
for an orthonormal $Q \in \R^{d \times k}$ whose column span defines a causal
subspace. Given a base input $x_b$ and source input $x_s$, the interchange
intervention replaces the in-subspace component:
\begin{equation}
h' = h_b - QQ^\top h_b + QQ^\top h_s .
\label{eq:intervention}
\end{equation}
IIA measures the fraction of interventions for which the model's output matches
the target counterfactual. High IIA indicates causal sufficiency: the subspace
contains enough information to drive behavior under interchange. But sufficiency
is not validity. A subspace may score high IIA because it encodes a memorized
lookup table, because correlated directions redundantly encode the same label, or
because the intervention distribution is too easy.
\citet{makelov2024illusion} demonstrated this concretely: subspace activation
patching can produce interpretability illusions where dormant pathways inflate IIA
without reflecting genuine causal structure. Our central claim is that IIA must be
paired with geometric diagnostics that test whether the subspace is a genuine
linear causal variable rather than an artifact of the linear function class.

\subsection{The Grassmannian manifold}
\label{sec:grassmannian_background}

The Grassmannian $\Gr(k,d)$ is the set of all $k$-dimensional linear subspaces
of $\R^d$. A DAS solution $Q \in \R^{d \times k}$ represents a point
$[Q] = \Span(Q)$ on this manifold; any $QR$ for $R \in O(k)$ represents the
same point. The geodesic distance between two subspaces is determined by their
\emph{principal angles}: if $Q_1, Q_2$ have orthonormal columns, the singular
values of $Q_1^\top Q_2$ are $\cos \theta_i$, and the geodesic distance is
\begin{equation}
\dgr([Q_1],[Q_2]) = \left(\sum_i \theta_i^2\right)^{1/2}
\label{eq:geodesic}
\end{equation}
\citep{edelman1998geometry, absil2008optimization}. Every application of DAS is
implicitly a search on this manifold. The question ``does DAS work?'' is
really ``does the true causal variable live on $\Gr(k,d)$?'' When it does, DAS
recovers a meaningful point. When it does not, DAS still returns a point on
$\Gr(k,d)$, but that point may be an artifact.

\subsection{Structured disentangled generative models}
\label{sec:structured_vae}

Semi-supervised deep generative models \citep{narayanaswamy2017learning} extend
the VAE framework \citep{kingma2014auto} by partitioning the latent space into
supervised and unsupervised factors. The generative model factors as
$p(x, y, z) = p(x \mid y, z) \, p(y) \, p(z)$, where $y$ is a labeled
(interpretable) variable and $z$ captures unlabeled nuisance variance. The
recognition model $q(y, z \mid x) = q(z \mid x, y) \, q(y \mid x)$ uses neural
network encoders, allowing nonlinear latent structure.

This framework is a natural generalization of DAS: DAS constrains the causal
variable to live on $\Gr(k,d)$ (a linear subspace), while the structured VAE
allows it to occupy a nonlinear manifold in activation space. The key advantage
is disentanglement: the supervised factor $y$ is trained to predict the
operation's output, while the unsupervised factor $z$ absorbs everything else.
If the causal variable is linear, the VAE encoder should learn a rotation
equivalent to DAS. If the causal variable is nonlinear (e.g., circular
structure from Fourier representations), the VAE can recover it where DAS
cannot.

The loss function combines the evidence lower bound (ELBO) with a supervised
classification term:
\begin{equation}
\mathcal{L} = \underbrace{-\mathbb{E}_{q}[\log p(x \mid y, z)]}_{\text{reconstruction}}
+ \beta \underbrace{D_\mathrm{KL}(q(y,z \mid x) \,\|\, p(y) p(z))}_{\text{regularization}}
+ \alpha \underbrace{\mathcal{L}_\text{CE}(y, \hat{y})}_{\text{supervision}}.
\label{eq:vae_loss}
\end{equation}

The encoder weights projecting to $y$ provide a linearized approximation of the
causal subspace; if this linearization has high overlap with the DAS subspace,
it confirms that the causal variable is approximately Grassmannian.

\subsection{Identifiability of the causal latent factor}
\label{sec:identifiability}

A central concern for any latent-variable method is \emph{identifiability}:
does the encoder recover the \emph{true} latent factors, or an arbitrary
reparameterization? The Structured pi-SAE satisfies the conditions of iVAE
\citep{khemakhem2020variational}: (i)~the task label $y$ is an observed
auxiliary variable; (ii)~the label-conditional prior
$p(z_\text{causal} \mid y) = \mathcal{N}(\mu_y, \sigma_y^2 I)$ is an
exponential family with per-label means that satisfy the rank condition;
(iii)~the ELBO reconstruction loss enforces an approximately injective
decoder.

\begin{proposition}[Causal variable recovery]
\label{prop:recovery}
Under the iVAE conditions (Assumption~\ref{ass:ivae} of
\citet{khemakhem2020variational}), the Structured pi-SAE encoder recovers
the true causal variable up to a component-wise bijection:
$\hat{z}_\text{causal} = \phi(z_\text{causal})$ where each $\phi_j$ depends
only on $z_j$. The proof follows directly from Theorem~1 of
\citet{khemakhem2020variational}; the supervised loss additionally resolves
the permutation ambiguity.
\label{ass:ivae}
\end{proposition}

\paragraph{Three consequences.}
\textbf{(1) NL-DAS violates condition~(iii)}:
its decoder is not reconstruction-constrained, so multiple latent codes may
map to the same activation and identifiability does not hold
(\S\ref{sec:nldas_vacuous}).
\textbf{(2) The component-wise ambiguity is harmless for IIA}:
interchange swaps $z_\text{causal}$ wholesale, so $\phi$ cancels.
\textbf{(3) Linear DAS is the special case $\phi \in O(k)$.}
When the true causal variable is linear, the structured encoder degenerates
to a rotation, matching the DAS solution.

\subsection{Grokking as a phase transition}

Grokking is a training regime in which models memorize the training set long
before generalizing \citep{power2022grokking}. Modular arithmetic models trained
with weight decay exhibit grokking reliably for some operations but not others
\citep{nanda2023progress, zhong2024clock}. Prior work analyzed grokking through
Fourier representations and circular structure in the embedding space. Our
perspective is complementary: we ask how the causal subspace recovered by DAS
changes across the memorization-to-generalization transition, and whether that
transition is deterministic or stochastic across random seeds.
\citet{liu2023grokking} frame grokking as compression---moving from a memorized
high-complexity solution to a generalizing low-complexity one; our Grassmannian
perspective gives this compression a geometric interpretation as convergence to a
specific point on $\Gr(k,d)$.


\section{Methods}
\label{sec:methods}

\subsection{Operations and models}

We study 14 binary operations over $\Z_p$ (Table~\ref{tab:operations}), chosen
to span group actions, polynomials, piecewise functions, and non-algebraic maps.
For each operation, we train a single-layer transformer with
$d_\text{model} = 128$, four attention heads ($d_\text{head} = 32$), and
$d_\text{MLP} = 512$ on examples of the form $(a, b, {=}) \mapsto y$. Training
uses the standard grokking setup: 30\% train split, AdamW with weight decay 1.0,
and 15{,}000--80{,}000 epochs depending on the operation. We classify a model as
\emph{grokked} if its final test cross-entropy is below 0.1.

\begin{table}[t]
\centering
\small
\begin{tabular}{llll}
\toprule
Operation & Formula & Category & Seeds tested \\
\midrule
Multiplication & $ab \bmod 113$ & Group action & 42, 2024 \\
Comp.\ addition & $(a{+}b) \bmod 91$ & Group action & orig, 42 \\
Subtraction & $(a{-}b) \bmod 113$ & Group action & 42 \\
Division & $a/b \bmod 113$ & Group action & 42 \\
Bitwise XOR & $a \oplus b \bmod 113$ & Group action & 42, 137, 2024 \\
\midrule
Sum of squares & $(a^2{+}b^2) \bmod 113$ & Polynomial & 42 \\
Cubic sum & $(a^3{+}b^3) \bmod 113$ & Polynomial & 42 \\
Cubing & $a^3 \bmod 113$ & Polynomial & 42 \\
Squaring & $a^2 \bmod 113$ & Polynomial & 42 \\
Polynomial & $(a^2{+}b) \bmod 113$ & Polynomial & 42 \\
Affine & $(2a{+}3b{+}5) \bmod 113$ & Polynomial & 42 \\
\midrule
Max & $\max(a,b) \bmod 113$ & Piecewise & 42 \\
Abs.\ difference & $|a{-}b| \bmod 113$ & Piecewise & 42 \\
Power & $a^b \bmod 113$ & Non-algebraic & 42, 137 \\
\bottomrule
\end{tabular}
\caption{Operation atlas. Group actions provide clean equivariance targets.
Polynomials test whether nonlinear maps admit linear causal variables. Piecewise
and non-algebraic operations probe the boundary.}
\label{tab:operations}
\end{table}

\subsection{DAS fitting and $k$-sweeps}

For each trained model, we cache residual-stream activations and train a DAS
rotation $Q \in \R^{d \times k}$ using interchange interventions for 400
gradient steps. We sweep $k \in \{2, 4, 6, 8, 10, 12, 16, 20, 24, 32\}$ to
estimate the intrinsic dimension $k^*$, defined as the smallest $k$ achieving
$\IIA \geq 0.95$. Each $k$ value is an independent DAS optimization, not a
nested subspace---$k$-sweep profiles reveal whether the causal variable is
inherently low-dimensional or requires many dimensions.

\subsection{Hard-example selection for IIA}
\label{sec:hard_examples}

Standard IIA evaluation includes many ``easy'' examples where the intervention
does not change the model's top-1 prediction, inflating IIA regardless of whether
the subspace is causally valid. We define a \emph{hard example} as a
(base, source) pair where: (1) the base and source have different correct outputs
$y_b \neq y_s$, (2) the model's prediction matches the correct output for both
base and source separately, and (3) the interchange intervention at the DAS
subspace \emph{can} flip the model's prediction from $y_b$ to $y_s$. Hard IIA
measures the fraction of hard examples where the flip actually occurs. This
selection follows the methodology of CPCA-init DAS
(Tower et al., forthcoming), which showed that standard IIA saturates too easily
while hard IIA reveals genuine causal structure.

\subsection{Equivariance testing}

For operations with a natural group action, we test whether the DAS subspace
respects the algebraic symmetry. Given input $(a, b)$ with DAS projection
$z = Q^\top h(a,b)$, we construct a shifted input $(a{+}1 \bmod p, b)$ and
check whether the DAS projection rotates by $2\pi/p$ radians. The equivariant
fraction is the proportion of test inputs satisfying this rotation property.
Random orthogonal subspaces of the same dimension serve as controls; across all
operations, random-subspace equivariance is below 1\%.

\subsection{Structured VAE for nonlinear causal variables}
\label{sec:vae_method}

For each trained grokking model, we cache residual-stream activations at the
post-MLP hook and train a structured VAE with:
\begin{itemize}[leftmargin=*]
\item $z_\text{causal}$: $k$-dimensional latent factor ($k \in \{2,4,8\}$),
  semi-supervised with the operation's output label via a classification head.
\item $z_\text{nuisance}$: $m$-dimensional unsupervised factor ($m = 15$) to
  absorb non-causal variance.
\item Encoder: 2-layer MLP ($d_\text{model} \to 128 \to z$) with ReLU.
\item Decoder: 2-layer MLP ($z \to 128 \to d_\text{model}$) with ReLU.
\end{itemize}

Training uses the combined ELBO + classification loss (Eq.~\ref{eq:vae_loss})
for 300 epochs with $\beta = 1.0$ and $\alpha = 10.0$.

\paragraph{VAE-based IIA.}
We evaluate the VAE's causal subspace by performing interchange interventions in
the \emph{latent} space: given base activation $h_b$ and source activation
$h_s$, we encode both, swap $z_\text{causal}$ while keeping $z_\text{nuisance}$
from the base, decode back to activation space, and measure whether the model
produces the source's output. This is the nonlinear analogue of
Eq.~\ref{eq:intervention}: instead of $QQ^\top$ projection, we use the
encode-swap-decode path.

\paragraph{VAE equivariance.}
We test whether the VAE's $z_\text{causal}$ exhibits equivariance under additive
shifts, using the same protocol as for DAS (shift input $a \to a+1$, check
whether $z_\text{causal}$ rotates consistently). Because the VAE encoder is
nonlinear, it can represent circular structure directly in a 2D latent space.

\subsection{Nonlinear DAS baseline}
\label{sec:nldas}

To isolate the contribution of the VAE's generative structure from the benefit
of nonlinearity alone, we compare against an unconstrained nonlinear featurizer.
Given MLP encoder $f: \R^d \to \R^d$ and decoder $g: \R^d \to \R^d$, the
nonlinear DAS (NL-DAS) intervention is:
\begin{equation}
h' = g\bigl(f(h_b) - QQ^\top f(h_b) + QQ^\top f(h_s)\bigr)
\label{eq:nldas}
\end{equation}
where $Q \in \R^{d \times k}$ is learned jointly with $f$ and $g$ by maximizing
interchange accuracy. Both $f$ and $g$ use the same architecture as the VAE
encoder (2-layer MLP, 256 hidden units, ReLU), matching the VAE's representational
capacity. Crucially, NL-DAS has \emph{no reconstruction constraint}: $f$ and $g$
are trained exclusively on interchange loss and need not be (approximately)
inverse to each other.

We also evaluate NL-DAS+recon, which adds a reconstruction penalty
$\lambda \| g(f(h)) - h \|^2$ to the training loss, forcing $f$ and $g$ to
approximate mutual inverses. This tests whether the unconstrained NL-DAS result
degrades when the encoder-decoder cannot be degenerate.

\subsection{Intervention faithfulness metrics}
\label{sec:faithfulness_metrics}

Binary IIA measures \emph{interchange effectiveness}: does swapping the
subspace component change the model's output to match the source? But
effectiveness alone does not imply that the intervention is \emph{faithful}---that
it modifies only the targeted causal variable while preserving all other
information. A method that achieves IIA~=~1.0 by overwriting the entire
activation with a class-specific template is effective but not faithful.

We introduce four complementary metrics to evaluate intervention faithfulness:

\paragraph{Continuous distributional metrics.}
Binary IIA saturates at 1.0 whenever the top-1 prediction flips, hiding
differences in intervention quality. We measure the full distributional impact:
\begin{itemize}[leftmargin=*]
\item \textbf{KL divergence}: $\KL(p_\text{clean} \| p_\text{iv})$, where
  $p_\text{clean}$ and $p_\text{iv}$ are the softmax distributions before and
  after intervention. Lower KL means the intervention preserves more of the
  clean distribution beyond the swapped variable.
\item \textbf{Jensen-Shannon divergence}: the symmetric alternative
  $\JS = \frac{1}{2}\KL(p \| m) + \frac{1}{2}\KL(q \| m)$ where $m = (p+q)/2$.
  Bounded in $[0, \ln 2]$, more stable than KL when distributions have disjoint
  support.
\item \textbf{Probability difference}: $p_\text{iv}(y_s) - p_\text{iv}(y_b)$,
  a continuous analog of binary IIA that captures the model's confidence margin,
  not just top-1 agreement.
\item \textbf{Normalized logit difference}:
  $\frac{l_\text{iv}(y_s) - l_\text{iv}(y_b)}{|l_\text{clean}(y_s) - l_\text{clean}(y_b)|}$,
  where $l$ denotes logits. A value of 1.0 means the intervention reproduces the
  clean model's margin exactly; values $\gg 1$ or $\ll 0$ indicate distortion.
\end{itemize}

\paragraph{Intervention diversity.}
For each source label $y_s$, we compute the standard deviation of the intervened
activations $h'$ across all base inputs with the same source. If the decoder
produces the same $h'$ regardless of which base input was used---varying only
with $y_s$---the intervention is a lookup table: it does not preserve
base-specific nuisance information. We define the \emph{diversity ratio}:
\begin{equation}
\rho = \frac{\mathbb{E}_{y_s}\bigl[\text{std}(h'_{y_s})\bigr]}
            {\mathbb{E}_{y_s}\bigl[\text{std}(h_{y_s})\bigr]}
\label{eq:diversity_ratio}
\end{equation}
where the numerator is over intervened activations sharing source label $y_s$
and the denominator is over the corresponding natural source activations.
$\rho \approx 0$ indicates a lookup table (all intervened activations for the
same $y_s$ collapse to a single point). $\rho \approx 1$ indicates the
intervention preserves the natural variation within each class.

\paragraph{Reconstruction fidelity.}
For encoder-decoder methods (NL-DAS, VAE), we measure $\|g(f(h)) - h\|^2$ on
held-out activations. A high reconstruction MSE with high IIA is a warning sign:
the decoder may be ``hallucinating'' the correct output rather than faithfully
reconstructing the activation with a modified causal component.
The generative constraint is not merely a regularizer---it is a necessary
condition for identifiability (Proposition~\ref{prop:recovery}), as shown
empirically by pi-VAE's failure despite having the correct prior direction
(Table~\ref{tab:ablation}).

\subsection{Cross-distribution generalization}
\label{sec:cross_distribution}

Standard train/test splits draw from the same distribution. For IOI on GPT-2, we
additionally evaluate under distribution shift: (1)~\textbf{disjoint names}---train
on one set of proper names, evaluate on a completely disjoint set; and
(2)~\textbf{disjoint templates}---train on sentence templates 1--3, evaluate on
templates 4--5 plus novel hard templates. A method that generalizes under
distribution shift is finding a genuine causal variable; one that fails is
exploiting distributional regularities in the training set.


\section{Results}
\label{sec:results}

\subsection{The atlas: a three-class partition}

The 14 operations partition cleanly into three classes based on whether DAS
recovers a Grassmannian causal variable (Table~\ref{tab:main_results}).

\begin{table}[t]
\centering
\small
\setlength{\tabcolsep}{3pt}
\begin{tabular}{lcccccl}
\toprule
Operation & Grok & Loss & IIA($k{=}2$) & $k^*$ & Equiv. & Class \\
\midrule
\multicolumn{7}{l}{\textit{Always Grassmannian}} \\
Multiplication & \checkmark & 0.000 & 1.000 & 2 & 0.995 & Always \\
Subtraction & \checkmark & 0.000 & 1.000 & 2 & 1.000 & Always \\
Division & \checkmark & 0.000 & 1.000 & 2 & 1.000 & Always \\
Bitwise XOR & \checkmark & 0.003 & 1.000 & 2 & 1.000 & Always \\
Cubic sum & \checkmark & 0.000 & 1.000 & 2 & 1.000 & Always \\
Sum of squares & \checkmark & 0.005 & 1.000 & 2 & 0.987 & Always \\
Max & \checkmark & 0.074 & 0.960 & 2 & 0.957 & Always \\
\midrule
\multicolumn{7}{l}{\textit{Stochastic}} \\
Comp.\ addition$^*$ & \checkmark & 0.000 & --- & --- & 0.998 & Stochastic \\
Comp.\ addition (s42) & $\times$ & 10.207 & 0.800 & 6 & 0.110 & Stochastic \\
Power (s137) & \checkmark & 0.088 & 0.940 & 4 & 0.965 & Stochastic \\
Power (s42) & $\times$ & 1.143 & 0.980 & 2 & 0.665 & Stochastic \\
\midrule
\multicolumn{7}{l}{\textit{Never Grassmannian}} \\
Cubing & $\times$ & 9.765 & 0.857 & 4 & 0.509 & Never \\
Squaring & $\times$ & 7.808 & 0.857 & 4 & 0.314 & Never \\
Abs.\ difference & $\times$ & 1.835 & 0.800 & $>$32 & 0.189 & Never \\
Polynomial & $\times$ & 21.550 & 0.720 & $>$32 & 0.015 & Never \\
Affine & $\times$ & 26.209 & 0.780 & $>$32 & 0.026 & Never \\
\bottomrule
\end{tabular}
\caption{Atlas results for 14 modular arithmetic operations. $k^*$ is the
smallest $k$ achieving $\IIA \geq 0.95$. The ``Always'' class contains operations
that grok and produce equivariance $>95\%$ across all tested seeds.
$^*$Original seed (unknown value); k-sweep not recorded.}
\label{tab:main_results}
\end{table}

\subsection{Linear DAS fails on grokked modular addition}
\label{sec:linear_das_fails}

The starkest evidence that the Grassmannian assumption can fail even for correct,
generalizing models comes from modular addition. We train DAS at
$k \in \{2, 4, 8, 16\}$ on residual-stream activations from a fully grokked
model (test loss $< 10^{-7}$, test accuracy 100\%). At \emph{every} subspace
dimension, DAS returns $\IIA = 0.0$---not low, but exactly zero.

This is not a failure of DAS optimization. The model has learned the correct
modular addition algorithm via Fourier representations \citep{nanda2023progress}:
the ``identity'' of each number $a$ is encoded as
$(\cos 2\pi f a / p, \sin 2\pi f a / p)$ for key frequencies $f$. This
representation lives on $S^1$ (a circle embedded in $\R^d$), not on a linear
subspace of $\R^d$. DAS searches $\Gr(k,d)$ for a flat plane that captures a
curved manifold---a category error, not an optimization failure.

This motivates the Structured pi-SAE approach (\S\ref{sec:vae_results}):
a nonlinear sparse encoder can map the circular representation to a structured
latent space where interchange interventions succeed. On grokked addition, the
Structured pi-SAE achieves IIA~=~1.0 at $k{=}2$, confirming that the causal
variable exists but is nonlinear.

\subsection{Stochastic grokking: the same operation, opposite outcomes}
\label{sec:stochastic}

The stochastic class provides the strongest evidence for our central claim.
Power ($a^b \bmod 113$) grokked at seed 137 (test loss 0.088, equivariance
96.5\%) but not at seed 42 (test loss 1.14, equivariance 66.5\%). Composite
addition ($a + b \bmod 91$) grokked at one seed (test loss 0.000, equivariance
99.8\%) but not at seed 42 (test loss 10.2, equivariance 11.0\%). Same
operation, same architecture, same hyperparameters---opposite outcomes from
random initialization alone. Multi-seed stability experiments (10 seeds per
stochastic operation) are underway and will be reported in the final version.

\subsection{Memorization artifacts: high IIA, no structure}
\label{sec:memorization}

Squaring and cubing achieve $\IIA = 0.86$ at $k = 2$ and $\IIA = 1.0$ by $k = 4$
without grokking (test losses 7.8 and 9.8). Equivariance exposes the artifact:
31.4\% for squaring and 50.9\% for cubing, far below the $>$95\% threshold. The
lesson: \textbf{high IIA with low-$k$ saturation is compatible with a lookup
table}. ``Works for interchange'' and ``is a Grassmannian causal variable'' are
different claims.

\subsection{$k$-sweep profiles correlate with equivariance}

Grassmannian variables saturate at $\IIA \approx 1.0$ by $k = 4$: their causal
structure is inherently 2-dimensional. Non-Grassmannian variables show gradual IIA
growth that may not reach 1.0 even at $k = 32$. Polynomial reaches
$\IIA = 0.72$ at $k = 2$ and requires $k \geq 8$ to approach 0.90.

\subsection{Structured pi-SAE recovers nonlinear causal variables}
\label{sec:vae_results}

We upgrade the vanilla structured VAE to a \textbf{Structured pi-SAE}: a
label-conditional sparse autoencoder that partitions the latent space into
$z_\text{causal}$ (supervised, $L_1$-penalized) and $z_\text{nuisance}$
(unsupervised). The ``pi'' denotes the label-conditional prior from
\citet{narayanaswamy2017learning}; the ``SAE'' denotes $L_1$ sparsity on
$z_\text{causal}$ with an expansion factor of $8\times$ (so $k$ causal
dimensions expand to $8k$ sparse features). This combines the identifiability
guarantees of auxiliary supervision \citep{khemakhem2020variational} with the
interpretability benefits of sparse dictionary learning
\citep{bricken2023towards}.

\paragraph{The $2 \times 2$ ablation.}
Both components---structured prior and sparsity---are necessary
(Table~\ref{tab:ablation}). On grokking tasks: (1)~the plain SAE (no
label-conditional prior) achieves IIA~=~1.0 on addition but degrades on harder
operations (multiplication~0.72, quartic sum~0.32); (2)~the pi-VAE (structured
prior, no sparsity) achieves IIA~$\approx$~0 across all operations;
(3)~\emph{only} the Structured pi-SAE achieves IIA~=~1.0 uniformly across all
grokked operations and IIA~$\approx$~0 on all non-grokked operations.

\begin{table}[t]
\centering
\small
\setlength{\tabcolsep}{4pt}
\begin{tabular}{@{}l|cc|cc@{}}
\toprule
& \multicolumn{2}{c|}{\textbf{Plain prior} $\mathcal{N}(0,I)$}
& \multicolumn{2}{c}{\textbf{Structured prior} $p(z \mid y)$} \\
\textbf{Operation} & VAE & SAE & pi-VAE & \textbf{Str.\ pi-SAE} \\
\midrule
\multicolumn{5}{l}{\textit{Grokked}} \\
Addition       & 0.02 & 1.00 & 0.00 & \textbf{1.00} \\
Multiplication & 0.00 & 0.72 & 0.00 & \textbf{1.00} \\
Quartic sum    & 0.02 & 0.32 & 0.00 & \textbf{1.00} \\
IOI (GPT-2)    & 0.56 & 0.83 & 0.95 & \textbf{0.98} \\
\midrule
\multicolumn{5}{l}{\textit{Non-grokked (correct answer: $\approx 0$)}} \\
Mixed product  & 0.03 & 0.01 & 0.01 & 0.04 \\
Squaring       & 0.00 & 0.00 & 0.00 & 0.00 \\
\bottomrule
\end{tabular}
\caption{$2{\times}2$ ablation: IIA at $k{=}2$.
Plain SAE works on addition but degrades on harder operations (multiplication
$0.72$, quartic sum $0.32$) because without the structured prior, the sparse
encoder has no target direction.  pi-VAE (structured prior, no sparsity)
achieves $0.00$ on grokked operations---the prior provides the right
direction but without sparsity the encoder spreads the causal signal across
all dimensions, which cannot be cleanly swapped.  Only Structured pi-SAE
achieves $1.00$ uniformly on all grokked operations and $\leq 0.04$ on all
non-grokked operations.}
\label{tab:ablation}
\end{table}

\paragraph{Reconstruction MSE as a falsifiability criterion.}
For non-grokked operations, reconstruction MSE is $11$--$24$; for grokked
operations, $0.6$--$1.2$.  High MSE with near-zero IIA means the model has
\emph{no structured causal variable to recover}---not that the method
failed.  This diagnostic is unavailable for NL-DAS, which returns IIA $=
0.6$--$0.8$ on non-grokked operations (false positives) because it has no
reconstruction constraint.

\paragraph{End-to-end intervention training.}
For language model tasks where the causal structure is more complex, we add an
\textbf{end-to-end (E2E) intervention CE loss}: after swapping $z_\text{causal}$
and decoding, we run the intervened activation through the remaining layers and
compute $-\log p(y_s \mid h')$. This differentiable objective directly
optimizes for intervention success, analogous to how DAS trains on interchange
accuracy. Training proceeds in two phases: 200 epochs of reconstruction +
classification warmup, followed by E2E intervention training with an additive
intervention $h' = h_b + \text{dec}(z_\text{swap}) - \text{dec}(z_\text{orig})$
that cancels reconstruction error.

\paragraph{Results on GPT-2 language tasks.}
We evaluate on five tasks from the MIB benchmark \citep{prakash2024mib} at
layer~8 of GPT-2-small: gender bias, hypernymy, greater-than, SVA, and
capitals. Table~\ref{tab:nlp_results} reports the headline result.

\begin{table}[t]
\centering
\small
\setlength{\tabcolsep}{4pt}
\begin{tabular}{lcccc}
\toprule
\textbf{Method} & \multicolumn{2}{c}{\textbf{Hypernymy}} & \multicolumn{2}{c}{\textbf{Gender bias}} \\
& $k{=}1$ & $k{=}4$ & $k{=}1$ & $k{=}4$ \\
\midrule
DAS              & 0.07 & 0.23 & 0.62 & 0.88 \\
NL-DAS           & 0.72 & 0.92 & 1.00 & 1.00 \\
Str.\ pi-SAE     & 0.58 & 0.52 & 0.41 & 0.52 \\
\textbf{Str.\ pi-SAE E2E} & \textbf{0.97} & \textbf{0.98} & \textit{pending} & \textit{pending} \\
\bottomrule
\end{tabular}
\caption{Standard IIA on GPT-2 language tasks (layer~8, additive intervention).
Structured pi-SAE with E2E training achieves 0.97 standard IIA on hypernymy at
$k{=}1$, where DAS gets 0.07 and even NL-DAS gets only 0.72. Strict IIA
(argmax = source token) is 0.90 at $k{=}1$.}
\label{tab:nlp_results}
\end{table}

On hypernymy, the Structured pi-SAE with E2E training achieves standard
IIA~=~0.97 at $k{=}1$---$14\times$ better than DAS (0.07) and $1.35\times$
better than NL-DAS (0.72). Under the stricter argmax criterion (source token
must beat all 50{,}000 vocabulary tokens), IIA remains 0.90. The E2E objective
closes the gap between training and evaluation by directly optimizing the
quantity we measure.

On the Indirect Object Identification task \citep{wang2022interpretability},
DAS achieves IIA~=~0.30 at $k{=}2$, consistent with the IOI causal
variable spanning multiple attention heads.  Structured pi-SAE achieves
$0.98$.  NL-DAS achieves $1.00$, but the diversity ratio $\rho \approx 0$
(vs.\ $\rho \approx 0.91$ for the pi-SAE) exposes it as vacuous
(Table~\ref{tab:nldas_vacuous}).  The identifiability guarantee
(Proposition~\ref{prop:recovery}) applies: each IOI sentence has a distinct
indirect object label, the per-label means $\mu_y$ are well-separated, and
the ELBO reconstruction constraint forces approximate decoder injectivity.

\subsection{Unconstrained nonlinear DAS is vacuous}
\label{sec:nldas_vacuous}

A natural response to DAS's linearity constraint is to prepend an MLP
featurizer: learn a nonlinear coordinate system in which the causal variable
\emph{becomes} linear, then apply standard DAS. This approach---which we call
NL-DAS (\S\ref{sec:nldas})---achieves IIA~=~1.000 on IOI at $k = 1$, compared
to DAS's 0.193. This looks like a dramatic improvement. It is not.

\paragraph{The lookup-table failure mode.}
NL-DAS is trained end-to-end on interchange loss with no reconstruction
constraint. The encoder $f$ and decoder $g$ need not be approximate inverses,
and the training objective provides no incentive for them to be. This creates a
degenerate solution: $f$ maps every activation to a space where the first
coordinate encodes the class label and the remaining coordinates are arbitrary;
$g$ ignores the base-specific content entirely and produces a ``canonical
activation for class $y_s$'' regardless of which base input was used.

The intervention diversity ratio (\S\ref{sec:faithfulness_metrics}) exposes this
failure. Table~\ref{tab:nldas_vacuous} reports results for IOI and three
grokking operations across $k \in \{1, 2, 4\}$.

\begin{table}[t]
\centering
\small
\setlength{\tabcolsep}{3pt}
\begin{tabular}{lc rrrr rr}
\toprule
Task & $k$ & DAS & NL-DAS & NL-DAS+r & VAE & $\rho_\text{NL}$ & $\rho_\text{VAE}$ \\
\midrule
IOI & 1 & 0.193 & 1.000 & \textbf{---} & 0.990 & \textbf{---} & \textbf{---} \\
IOI & 2 & 0.297 & 1.000 & \textbf{---} & 0.980 & \textbf{---} & \textbf{---} \\
IOI & 4 & 0.556 & 1.000 & \textbf{---} & 0.978 & \textbf{---} & \textbf{---} \\
Addition & 1 & 0.006 & 0.194 & \textbf{---} & 0.000 & \textbf{---} & \textbf{---} \\
Addition & 4 & 0.289 & 0.883 & \textbf{---} & 0.089 & \textbf{---} & \textbf{---} \\
Mult.\ & 4 & 0.344 & 0.961 & \textbf{---} & 0.000 & \textbf{---} & \textbf{---} \\
Squaring & $*$ & 0.000 & 0.011 & \textbf{---} & 0.000 & \textbf{---} & \textbf{---} \\
\bottomrule
\end{tabular}
\caption{NL-DAS achieves high IIA by learning degenerate encoder-decoders.
  The VAE column reports Structured pi-SAE results.
  $\rho$ is the intervention diversity ratio (Eq.~\ref{eq:diversity_ratio});
  NL-DAS+r adds a reconstruction penalty.
  \textbf{---} marks values pending from hard-mode experiments.
  All controls (random labels, reconstruction-only, untrained, unconstrained
  plain VAE) achieve IIA~$< 0.09$, confirming the Structured pi-SAE is not
  cheating.}
\label{tab:nldas_vacuous}
\end{table}

Three lines of evidence confirm NL-DAS is vacuous:

\begin{enumerate}[leftmargin=*]
\item \textbf{Controls rule out the structured VAE cheating.} Random-label VAE
  (IIA $\leq 0.02$), reconstruction-only VAE with $\alpha = 0$
  (IIA $\leq 0.005$), untrained VAE (IIA = 0.000), and an unconstrained plain
  VAE with no causal/nuisance split (IIA $\leq 0.08$) all fail. The structured
  VAE's success requires correct supervision \emph{and} the
  causal-nuisance partition. This is the empirical counterpart of Proposition~\ref{prop:recovery}: NL-DAS
  fails Assumption~\ref{ass:ivae}(iii) because its decoder is not injective,
  so the iVAE identifiability guarantee does not apply
  \citep{khemakhem2020variational}. Without inductive biases, unsupervised
  disentanglement is impossible \citep{locatello2019challenging}; the
  structured prior and reconstruction constraint are not optional regularizers
  but necessary conditions.

\item \textbf{NL-DAS+recon degrades.} When forced to reconstruct its input
  ($g(f(h)) \approx h$), NL-DAS can no longer achieve degenerate solutions.
  Results with the reconstruction penalty are forthcoming.

\item \textbf{Cross-distribution generalization.}
  NL-DAS performance under distribution shift (disjoint names and novel
  templates), where the encoder cannot memorize specific activation patterns
  from training, will be reported in the final version.
\end{enumerate}

\paragraph{Why the VAE is not vacuous.}
The structured VAE's ELBO provides three constraints that prevent the
lookup-table failure mode: (1)~the reconstruction loss forces the decoder to
produce activations close to the original, not class-specific templates;
(2)~the KL regularization prevents the encoder from collapsing the latent space
to a few discrete points; and (3)~the causal-nuisance partition ensures that
$z_\text{nuisance}$ absorbs base-specific information, so the decoder must use
it. Together, these constraints mean that swapping $z_\text{causal}$ while
preserving $z_\text{nuisance}$ genuinely modifies only the causal component.

\paragraph{Implications.}
Any nonlinear method for causal abstraction must be evaluated with intervention
faithfulness metrics, not just IIA. The NL-DAS failure mode is not specific to
our architecture---it applies to any unconstrained encoder-decoder trained
exclusively on interchange loss. Prior work proposing neural network featurizers
for causal abstraction \citep{geiger2024causal} should be re-evaluated with
diversity ratio and reconstruction checks.

\subsection{Continuous metrics reveal what binary IIA hides}
\label{sec:continuous_metrics}

Binary IIA treats a 51\% probability shift the same as a 99\% probability shift.
For the always-Grassmannian operations where IIA saturates at 1.0, continuous
metrics reveal important quality differences between methods. Full continuous
metrics tables (KL, JS, probability difference, normalized logit difference)
comparing DAS, NL-DAS, and Structured pi-SAE across all operations will be
reported in the final version.

\subsection{Surprising positive and negative cases}

Cubic sum ($a^3 + b^3 \bmod 113$) achieves perfect equivariance (100\%) despite
cubing alone reaching only 50.9\%. The boundary is not ``group action versus
polynomial'' but whether the operation's binary structure provides sufficient
additive symmetry.

Affine ($2a + 3b + 5 \bmod 113$) is a linear function that \emph{fails} to
produce a Grassmannian variable. Under an additive shift $a \mapsto a + g$, the
output shifts by $2g$, not $g$. Equivariance requires the operation to be a group
action, not merely a linear function.

Max ($\max(a,b) \bmod 113$) achieves 95.7\% equivariance despite not being a
group action. It is piecewise equivariant: when the argmax is stable under the
shift, max behaves like a shifted identity.


\section{The Mechanism: How Grokking Creates Causal Structure}
\label{sec:mechanism}

The atlas (\S\ref{sec:results}) established \emph{what}: a sharp three-class
partition governed by grokking. This section asks \emph{how} and \emph{why}.
We trace the emergence of causal structure through five lenses: the minimal
architecture required, the fine-grained dynamics of Fourier crystallization,
spectral compression during the phase transition, the relationship between
weight-space structure and activation-space causal variables, and the
gauge-freedom of the causal subspace itself.

\subsection{Minimal architecture for grokking}
\label{sec:dmf}

What is the simplest model that groks? We compare a depth-2 deep matrix
factorization (DMF) $y = W_2 W_1 x$ against variants with nonlinearity
(DMF + ReLU), a linearized transformer (identity attention + ReLU MLP),
a softmax-only transformer (no MLP), and the full single-layer transformer.
All are trained on modular multiplication ($\Z_{113}$) with weight decay 1.0
for 50{,}000 epochs.

\begin{table}[t]
\centering
\small
\begin{tabular}{lrrl}
\toprule
Architecture & Grok epoch & Fourier align. & Groks? \\
\midrule
DMF + ReLU     & $\sim$2{,}000   & 0.095 & Yes (18$\times$ faster) \\
MLP-only       & $\sim$2{,}500   & 0.99  & Yes \\
Linearized TF  & $\sim$14{,}000  & 0.000 & Yes (no Fourier) \\
Full transformer & $\sim$37{,}000 & 0.10$\to$0.54 & Yes \\
Linear DMF     & ---          & 0.078 & \textbf{No} \\
Softmax-only   & ---          & 0.103 & \textbf{No} \\
\bottomrule
\end{tabular}
\caption{Grokking speed hierarchy. Nonlinearity is necessary (linear DMF never
groks despite nuclear norm dropping from 7{,}340 to 3{,}877). Attention is not
(DMF + ReLU groks 18$\times$ faster). Softmax alone is insufficient.}
\label{tab:grokking_speed}
\end{table}

Three findings emerge. First, \textbf{nonlinearity is necessary}: the linear
DMF never generalizes despite 50{,}000 epochs. Its nuclear norm drops from
7{,}340 to 3{,}877---spectral compression occurs---but test accuracy remains
$\sim$0\%. Compression alone is not sufficient for grokking; the nonlinear
activation function that enables the Fourier representation is essential.

Second, \textbf{attention is not necessary for grokking on modular
multiplication}: DMF + ReLU groks 18$\times$ faster than the full
transformer. The softmax attention mechanism adds enormous overhead to the
grokking process without being required for the underlying algorithm on this
operation. Whether this extends to other operations or multi-task settings
remains open.

Third, the \textbf{linearized transformer groks without Fourier structure}:
it achieves 100\% test accuracy by epoch 14{,}000 but its Fourier alignment
remains at 0.000 throughout. Grokking and Fourier representation are
dissociable---grokking can occur via non-Fourier algorithms. The Fourier
representation is one route to generalization, not the only one.

\subsection{Fourier structure emerges after generalization}
\label{sec:fourier_trajectory}

Prior work identified Fourier representations as the mechanism underlying
grokking in modular arithmetic \citep{nanda2023progress, zhong2024clock}.
A natural assumption is that Fourier structure \emph{drives} grokking:
the model discovers the Fourier basis, then generalizes. Our fine-grained
trajectory analysis (1{,}002 checkpoints at 100-epoch resolution) reveals
the opposite: \textbf{generalization precedes Fourier crystallization}.

\begin{table}[t]
\centering
\small
\begin{tabular}{rrl}
\toprule
Epoch & Test accuracy & Fourier alignment \\
\midrule
37{,}000 & 65.6\% & 0.070 \\
37{,}500 & 91.3\% & 0.093 \\
38{,}000 & 98.6\% & 0.143 \\
38{,}500 & 99.7\% & 0.260 \\
39{,}500 & 99.7\% & 0.498 \\
40{,}000 & 99.8\% & 0.543 \\
\bottomrule
\end{tabular}
\caption{Fine-grained Fourier trajectory for modular multiplication.
Test accuracy crosses 90\% at epoch 37{,}500; Fourier alignment crosses 0.30
approximately 1{,}000 epochs later. The model generalizes \emph{before} its
weights crystallize into the Fourier basis.}
\label{tab:fourier_lag}
\end{table}

Test accuracy crosses 90\% at epoch 37{,}500, but Fourier alignment does not
cross 0.30 until approximately epoch 38{,}500---a lag of $\sim$1{,}000
epochs. The model generalizes before its weights fully organize into the
Fourier basis.

\paragraph{The AGOP double-peak.}
Tracking the Average Gradient Outer Product (AGOP) Fourier alignment through
training reveals an unexpected non-monotonic pattern. For the full transformer:

\begin{table}[t]
\centering
\small
\begin{tabular}{rrrrr}
\toprule
Epoch & Test acc. & AGOP FA & Weight FA & AGOP erank \\
\midrule
5{,}000   & 0.4\%   & 0.108 & 0.080 & 10.0 \\
25{,}000  & 6.3\%   & 0.712 & 0.081 & 9.8 \\
30{,}000  & 19.1\%  & \textbf{0.894} & 0.075 & 8.7 \\
35{,}000  & 95.8\%  & 0.662 (dip) & 0.270 & 5.4 \\
40{,}000  & 100\%   & \textbf{0.974} & 0.864 & 4.7 \\
45{,}000  & 100\%   & 0.888 & 0.977 & 4.4 \\
\bottomrule
\end{tabular}
\caption{AGOP trajectory for modular multiplication. The Fourier alignment of
the gradient outer product shows a double-peak: 0.894 at epoch 30{,}000,
followed by a dip to 0.662 during the grokking transition (epoch 35{,}000),
then crystallization at 0.974 (epoch 40{,}000). The \emph{gradient}
discovers Fourier structure before the \emph{weights} do (weight FA is
0.075 when AGOP FA is 0.894).}
\label{tab:agop}
\end{table}

The AGOP alignment peaks at 0.894, then \emph{dips} to 0.662 during the
actual grokking transition (epoch 35{,}000), before crystallizing at 0.974
post-grokking. The dip coincides with the period of rapid generalization---the
phase transition temporarily disrupts the gradient's Fourier structure before
it stabilizes. By contrast, the DMF's AGOP alignment remains flat at
0.06--0.08 throughout 50{,}000 epochs, confirming that the linear model never
discovers Fourier structure even in its gradients.

The key implication: the gradient ``knows'' the Fourier basis $\sim$10{,}000
epochs before the weights do (AGOP FA is 0.894 at epoch 30{,}000 while weight
FA is still 0.075). Grokking may be the process by which weight-space catches
up to gradient-space structure.

\paragraph{Fourier structure lives in attention, not MLPs.}
In a grokked 1-layer model ($p = 113$), we measure Fourier selectivity
across all components: 0 of 512 MLP neurons are Fourier-selective, and the
embedding has a nearly flat Fourier spectrum (participation ratio = 3.93).
The Fourier structure resides entirely in attention: the QK circuit selects
individual frequency pairs (rank-1 structure), while the OV circuit operates
in a broader $k$-dimensional subspace of Fourier modes. This decomposition
is consistent with the ``clock'' interpretation of \citet{zhong2024clock}.

\subsection{Spectral compression during grokking}
\label{sec:spectral_compression}

The factored feedforward product matrix $W_\text{FF} = W_\text{in} W_\text{out}$
undergoes dramatic rank collapse during grokking. Tracking the effective rank
(exponential of the Shannon entropy of the normalized singular value spectrum)
through training:

\begin{center}
\small
\begin{tabular}{rrl}
\toprule
Epoch & FF effective rank & Stage \\
\midrule
34{,}000 & 48 & Pre-grokking \\
37{,}500 & 14 & Transition \\
41{,}000 & 9 & Post-grokking \\
\bottomrule
\end{tabular}
\end{center}

The 5.3$\times$ compression from effective rank 48 to 9 confirms that
grokking compresses the model's representational capacity into a
low-dimensional subspace. This compression is necessary but not sufficient:
the linear DMF also achieves spectral compression (nuclear norm
$7{,}340 \to 3{,}877$) without grokking (\S\ref{sec:dmf}). What distinguishes
successful grokking is that the compressed subspace acquires
\emph{structure}---specifically, equivariance under the operation's symmetry
group.

\subsection{Weight-space SVD predicts causal geometry}
\label{sec:svd_vs_das}

If grokking compresses the solution into a low-rank subspace, can we read the
causal variable directly from the weight matrices without fitting DAS? We
compare SVD of the feedforward product matrix against DAS fitted to
activations, measuring IIA for both.

\begin{table}[t]
\centering
\small
\begin{tabular}{llrl}
\toprule
Method & Site & Best $k$ & IIA \\
\midrule
SVD(FF) & attn\_out & 32 & \textbf{0.999} \\
DAS     & attn\_out & 8  & 0.150 \\
DAS     & mlp\_out  & 8  & 0.631 \\
SVD(FF) & mlp\_out  & 64 & 0.444 \\
\bottomrule
\end{tabular}
\caption{SVD of the feedforward product matrix vs.\ DAS on a grokked
multiplication model. At the attn\_out site, SVD achieves IIA = 0.999
\emph{with no optimization} (DAS: 0.150 after 1{,}000 gradient steps).
The Grassmannian distance between SVD and DAS subspaces is
1.5--8.2~radians (near-orthogonal).}
\label{tab:svd_vs_das}
\end{table}

At the attention output site, SVD(FF) achieves $\IIA = 0.999$ at $k = 32$
with \emph{zero optimization}---no gradient steps, no interchange
training---while DAS achieves only 0.150 at $k = 8$ after 1{,}000 gradient
steps. The comparison is not dimension-matched: SVD requires $4\times$ the
subspace dimension, reflecting that the top singular vectors of $W_\text{FF}$
include directions relevant to the causal variable but also spectral
directions that are not causally active. Nevertheless, the fact that a
zero-cost weight-space decomposition outperforms an optimized
activation-space method is striking.

The Grassmannian distance between the SVD and DAS subspaces ranges from
1.5 to 8.2 radians---they are near-orthogonal. This is consistent with
the gauge-freedom finding (\S\ref{sec:gauge_freedom}): many different
subspaces carry the same causal information. The SVD subspace is one
valid parameterization; the DAS subspace is another. Neither is
``correct''---both are equivalent points in the orbit of the gauge
group acting on $\Gr(k,d)$.

The site dependence is notable: SVD(FF) is most informative at attn\_out
(where the attention output has been computed) but weaker at mlp\_out
(where the MLP has further mixed the representation). DAS shows the
opposite pattern, achieving 0.631 at mlp\_out. This suggests that
weight-space and activation-space methods have complementary strengths:
SVD captures the structural potential of the weight matrices, while DAS
captures the functional use of that structure under the data distribution.

\subsection{Subspace identity is gauge freedom}
\label{sec:gauge_freedom}

If grokking produces a specific causal subspace, we might expect different
random seeds to converge to the \emph{same} point on $\Gr(k,d)$. They do
not. Ten seeds of modular multiplication all grok successfully ($\IIA$
0.94--1.00), but the recovered DAS subspaces are nearly orthogonal: pairwise
overlap $\approx 0.008$, where random $k$-dimensional subspaces of
$\R^{128}$ would have overlap $k/d = 0.016$. The grokked subspaces are
\emph{less} aligned than chance.

\paragraph{The basin of attraction is enormous.}
Perturbing the DAS subspace by rotating it 1+ radians on $\Gr(k,d)$
barely drops IIA. The causal variable is not a fragile low-dimensional
needle in a high-dimensional haystack---it is a robust feature of the
representation that can be read out from exponentially many subspaces.

\paragraph{Spectral graph structure.}
Computing the pairwise Grassmannian distance between all 10 seeds'
subspaces yields a nearly uniform distance matrix:
$d_{\Gr} \in [1.93, 2.19]$. The subspaces are equidistant from each other,
forming a ``spectral simplex'' on the Grassmannian rather than clustering
around a preferred direction.

This finding has a concrete implication for mechanistic interpretability:
\textbf{the identity of a causal subspace is not a meaningful property}.
What matters is the equivariance of the learned function \emph{within} the
subspace, not the subspace itself. Two analyses of the same model at
different random seeds would find orthogonal DAS solutions---but both
solutions are equally valid because the model treats the choice of
subspace as gauge freedom. This parallels the rotational gauge freedom
in the factorized transformer's QK circuit, where $W_Q \to W_Q R$,
$W_K \to W_K R$ leaves $W_Q W_K^\top$ invariant.

\subsection{Structural robustness under un-training}
\label{sec:un_training}

If grokking creates causal structure, does un-training---continuing to train
on the wrong operation's labels---destroy it? We present preliminary results
($n{=}1$ per condition) from three grokked models retrained with shuffled
labels from a different operation.

\begin{table}[t]
\centering
\small
\begin{tabular}{lrrrl}
\toprule
Operation & Still grokked? & Test loss & OV erank & Interpretation \\
\midrule
Multiplication & No & 20.4 & 9.19 & Structure degraded \\
Comp.\ addition & \textbf{Yes} & 9e-6 & 5.46 & Structure preserved \\
Cubing & No & 9.19 & 6.23 & Higher erank \\
\bottomrule
\end{tabular}
\caption{Un-training results: retraining grokked models on wrong-operation
labels. Composite addition retains its grokked structure (test loss 9e-6)
even under wrong labels---the learned algorithm is robust to label
corruption. Multiplication loses grokking (erank increases from the grokked
baseline to 9.19).}
\label{tab:un_training}
\end{table}

Composite addition retains its grokked structure completely---test loss
remains $9 \times 10^{-6}$ and OV effective rank stays at 5.46. The learned
algorithm is so deeply embedded in the weight structure that training on
wrong labels cannot dislodge it. Multiplication, by contrast, loses
grokking: its erank increases to 9.19, and test
loss rises to 20.4.

The asymmetry is informative. Composite addition uses $p = 91$ (composite
modulus), while multiplication uses $p = 113$ (prime). The composite
modulus may create more redundant encoding of the Fourier structure,
making it harder to un-learn. This connects to the stochastic grokking
finding (\S\ref{sec:stochastic}): composite addition's grokking is
seed-dependent, but once it grokks, the structure is more robust than
prime-modulus operations.

\subsection{Stratum classification: toy models vs.\ pretrained}
\label{sec:strata}

We test whether the singular value structure of weight matrices---classifying
each head's OV and QK circuits into strata based on spectral gap, cumulative
variance, and effective rank---provides an alternative lens on causal
geometry.

In grokked 1-layer models ($p = 113$), the classification works: all four
heads have low-rank OV (erank 2.1--6.2) and rank-1 QK (gap ratio $> 3$),
consistent with the Fourier mechanism where QK selects frequency pairs and
OV operates in the active Fourier subspace.

In GPT-2-small, it fails completely. Of 144 heads, 143 fall into the
moderate-rank stratum ($\mathcal{M}_\Sigma$) regardless of function.
Circuit heads (IOI name movers, induction, s-inhibition) have \emph{higher}
average OV erank (54.6) than non-circuit heads (49.5)---the opposite of
what low-rank structure would predict. The explanation is superposition:
each head participates in multiple tasks, producing high effective rank
even if each task-specific contribution is low-rank.

The methodological lesson: \textbf{functional diagnostics (equivariance)
outperform structural diagnostics (spectral profile)} for classifying
causal geometry. Equivariance cleanly separates the three-class partition
without spectral analysis. Task-conditioned strata---projecting OV into
a task-specific DAS subspace before classifying---may rescue the approach
but remain untested.


\section{Discussion}
\label{sec:discussion}

\subsection{A hierarchy of causal validity}

Our results establish a four-level hierarchy for evaluating causal variable
claims, where each level adds a stronger requirement:

\begin{enumerate}
\item \textbf{Interchange effectiveness} (IIA; necessary, not sufficient): High
  IIA means the method supports interchange but does not distinguish genuine
  causal structure from memorization (squaring achieves $\IIA = 1.0$ at $k = 4$
  without grokking) or from degenerate encoder-decoders (NL-DAS achieves
  $\IIA = 1.0$ at $k = 1$ via lookup tables).

\item \textbf{Interchange faithfulness} (diversity ratio, reconstruction;
  necessary for nonlinear methods): A faithful intervention modifies the causal
  variable while preserving all other information. The diversity ratio
  $\rho$ (Eq.~\ref{eq:diversity_ratio}) and reconstruction MSE detect
  lookup-table failure modes. Without faithfulness, ``high IIA'' means only
  ``the decoder can produce the right output,'' not ``the encoder identified
  the right variable.'' Linear DAS is faithful by construction (the projection
  $QQ^\top$ preserves the orthogonal complement), but nonlinear methods must
  demonstrate faithfulness empirically.

\item \textbf{Distributional fidelity} (KL, JS, continuous metrics;
  distinguishes quality among effective methods): Among methods achieving
  IIA~$\approx$~1.0, continuous distributional metrics reveal how much of
  the clean distribution is preserved versus distorted. A method that flips
  the top-1 prediction (high IIA) while scrambling the rest of the distribution
  (high KL) is less valid than one that reproduces the full counterfactual
  distribution.

\item \textbf{Structural diagnostics} (equivariance, cross-distribution
  generalization): High equivariance means the recovered variable transforms
  coherently under the operation's symmetry group---the strongest evidence that
  the method has found a genuine causal variable rather than exploiting
  distributional artifacts. Cross-distribution generalization (disjoint names,
  novel templates) provides the analogous test for natural-language tasks that
  lack algebraic symmetries.
\end{enumerate}

The hierarchy has an important asymmetry. Linear DAS automatically satisfies
level~2 (faithfulness) because the orthogonal projection modifies only the
subspace component, but it may fail level~1 when the causal variable is
nonlinear. Unconstrained nonlinear methods easily satisfy level~1 but fail
level~2---they can achieve perfect IIA by learning degenerate solutions. The
Structured pi-SAE satisfies both because the ELBO provides the reconstruction and
regularization constraints that prevent degeneracy.

\subsection{Grokking enables causal structure---linear or nonlinear}

The stochastic grokking discovery shows that grokking is the process by which a
model's internal representation transitions from an unstructured memorization
encoding to a structured causal encoding. For operations with linear algebraic
structure, this structured encoding lives on $\Gr(k,d)$ and DAS recovers it.
For operations with nonlinear structure, the encoding lives on a curved manifold
and requires nonlinear methods to access.

The mechanism section (\S\ref{sec:mechanism}) reveals that this transition
involves three dissociable processes: spectral compression (effective rank
$48 \to 9$), Fourier crystallization (which \emph{lags} generalization by
$\sim$1{,}000 epochs), and the selection of one among exponentially many
equivalent causal subspaces (gauge freedom). The AGOP analysis
(\S\ref{sec:fourier_trajectory}) provides a concrete timeline: the gradient
discovers Fourier structure $\sim$10{,}000 epochs before the weights do,
and the actual phase transition is marked by a temporary \emph{dip} in
gradient-space Fourier alignment before crystallization.

The minimal architecture experiments (\S\ref{sec:dmf}) further clarify:
nonlinearity is necessary for grokking, but attention is not. DMF + ReLU
groks 18$\times$ faster than the full transformer, and the linearized
transformer groks without any Fourier structure at all. These findings
suggest that Fourier representations are the dominant grokking mechanism
for transformers with softmax attention, but grokking itself is a more
general phenomenon---the transition from memorization to a structured
causal encoding---that can occur through multiple algorithmic routes.

The key insight is that grokking is about \emph{causal structure}, not
\emph{linear} causal structure. The Grassmannian perspective captures the linear
case precisely, but the broader phenomenon is the emergence of any structured
causal variable---linear or nonlinear---during the generalization transition.

\subsection{Implications for language model interpretability}

For practitioners applying DAS or nonlinear causal abstraction to language
models:
\begin{enumerate}
\item \textbf{Never report IIA alone.} IIA is necessary but insufficient at
  every level: linear DAS on memorized models (squaring, $\IIA = 1.0$) and
  unconstrained nonlinear methods (NL-DAS, $\IIA = 1.0$) both achieve
  perfect IIA vacuously. Always report at least one faithfulness metric
  (diversity ratio for nonlinear methods, equivariance or cross-distribution
  generalization for any method).
\item \textbf{Use continuous distributional metrics.} KL divergence and
  normalized logit difference distinguish among methods that all achieve
  $\IIA \approx 1.0$. Binary IIA treats a 51\% probability shift the same as
  99\%; continuous metrics capture this difference.
\item \textbf{Be suspicious of unconstrained nonlinear featurizers.} Any
  encoder-decoder trained end-to-end on interchange loss without reconstruction
  or regularization constraints can learn a lookup table. The ELBO
  (or an equivalent generative constraint) is necessary to prevent this
  failure mode.
\item \textbf{Test cross-distribution generalization.} For natural-language
  tasks, evaluate on disjoint entity sets or novel templates. A method that
  fails under distribution shift is exploiting surface regularities, not
  recovering causal structure.
\item \textbf{Consider nonlinear methods} when DAS IIA is low. Low IIA does not
  mean the causal variable does not exist---it may mean DAS is searching the
  wrong function class. But validate with faithfulness metrics.
\end{enumerate}

\subsection{Weight-space vs.\ activation-space methods}

The SVD result (\S\ref{sec:svd_vs_das}) reveals a surprising complementarity.
Weight-space methods (SVD of $W_\text{in} W_\text{out}$) achieve near-perfect
IIA at attn\_out \emph{without any optimization}, outperforming 1{,}000 steps
of DAS gradient descent. This suggests that for grokking models, the causal
variable is directly legible in the weight matrices---DAS's optimization
is recovering structure that was already ``written down'' in the weights.

The practical implication is that weight-space analysis should precede
activation-space methods: if SVD of the relevant weight product achieves
high IIA, DAS optimization may be unnecessary. If SVD fails but DAS
succeeds, the discrepancy suggests the causal variable involves
computation beyond simple linear readout from the weight matrices.

\subsection{Limitations}

Our grokking experiments use single-layer transformers on synthetic tasks.
Whether the grokking--Grassmannian link holds in deeper models or natural
language is an open question, though the $k$-sweep and equivariance
diagnostics apply regardless. The Structured pi-SAE introduces a model-selection
problem (architecture, $\beta$, $\alpha$) that DAS avoids. Multi-seed
stability experiments (10 seeds per stochastic operation) are underway. The
connection between the Structured pi-SAE latent space and the Fourier
representation structure identified by \citet{nanda2023progress} deserves
further theoretical analysis. The stratum classification
(\S\ref{sec:strata}) fails on pretrained models; task-conditioned strata
remain untested. The un-training experiments (\S\ref{sec:un_training}) are
preliminary ($n{=}1$ per condition, three operations)---multi-seed
replication and broader operation coverage are needed before drawing firm
conclusions about un-training robustness.


\section{Related Work}
\label{sec:related}

\paragraph{Causal abstraction.}
Causal abstraction provides the theoretical foundation for DAS
\citep{geiger2021causal, geiger2023finding}. \citet{wu2024interpretability}
scaled DAS via Boundless DAS. \citet{makelov2024illusion} identified
interpretability illusions where dormant pathways inflate IIA. Our geometric
diagnostics address a complementary problem: even within the linear intervention
class, high IIA does not imply genuine linear structure.

\paragraph{Grokking.}
Grokking was introduced by \citet{power2022grokking} and mechanistically
analyzed through Fourier representations \citep{nanda2023progress,
zhong2024clock}. \citet{liu2023grokking} frame grokking as compression;
our spectral analysis (\S\ref{sec:spectral_compression}) gives this
compression concrete numbers (effective rank $48 \to 9$, $5.3\times$) and
shows it is necessary but not sufficient---linear DMFs compress without
grokking. Our AGOP analysis (\S\ref{sec:fourier_trajectory}) adds temporal
resolution absent from prior work: the gradient discovers Fourier structure
$\sim$10{,}000 epochs before the weights, and generalization precedes
Fourier crystallization by $\sim$1{,}000 epochs. Our contribution is to
show that grokking has a specific consequence for causal abstraction: it is
the process that converts non-Grassmannian representations into Grassmannian
ones (for linear operations) or into structured nonlinear ones (for nonlinear
operations). Stochastic grokking has been observed via the slingshot mechanism
\citep{thilak2022slingshot} but not connected to causal geometry.

\paragraph{Spectral analysis of neural networks.}
Random matrix theory has been applied to weight matrices for understanding
training dynamics \citep{martin2021implicit}. Our stratum classification
(\S\ref{sec:strata}) extends this to a functional taxonomy, but reveals an
important limitation: spectral profiling has discriminative power for
single-task models (where it correctly identifies QK as rank-1 and OV as
low-rank in grokked models) but loses all discriminative power under
superposition (143/144 GPT-2 heads are $\mathcal{M}_\Sigma$ regardless of
function). Task-conditioned spectral analysis may bridge this gap but remains
untested.

\paragraph{Disentangled representations and causal variables.}
\citet{narayanaswamy2017learning} introduced semi-supervised VAEs for
disentanglement. Recent work on identifiable representations
\citep{khemakhem2020variational} establishes conditions under which latent
factors can be recovered. \citet{locatello2019challenging} proved that
unsupervised disentanglement is impossible without inductive biases; our
structured prior is precisely the auxiliary supervision their theorem
requires. Our application differs in that we evaluate
disentanglement through causal interchange (IIA), not reconstruction quality,
connecting the VAE literature to causal abstraction.

\paragraph{Representation geometry.}
Grassmannian methods have been applied to subspace tracking and domain adaptation
\citep{edelman1998geometry, gopalan2011domain}. Linear representation hypotheses
in interpretability \citep{marks2023geometry, engels2024language,
park2024linear} argue that many concepts are encoded in linear subspaces; our
work tests this hypothesis causally rather than correlationally.

\paragraph{Nonlinear causal abstraction.}
Concurrent work on nonlinear interchange interventions
\citep{geiger2024causal} uses neural network featurizers before applying DAS.
Our results show that this approach is vulnerable to the lookup-table failure
mode (\S\ref{sec:nldas_vacuous}): unconstrained featurizers achieve perfect IIA
by learning degenerate encoder-decoders. Our Structured pi-SAE approach uses the
ELBO to provide the reconstruction and regularization constraints that prevent
this degeneracy, connecting to the identifiability literature
\citep{khemakhem2020variational, sutter2020multimodal} which establishes that
auxiliary supervision is necessary for recovering latent factors. The NL-DAS
failure mode and our faithfulness metrics provide empirical evidence for this
theoretical requirement in the causal abstraction setting.


\section{Conclusion}
\label{sec:conclusion}

DAS searches for causal variables in a specific function class: linear subspaces
on the Grassmannian. Our atlas of 14 operations maps where this search succeeds
and where it fails. The boundary is sharp, governed by grokking, and stochastic
for operations near the transition.

Three practical implications follow. First, IIA should never be reported alone:
$\IIA = 1.000$ is compatible with a genuine causal variable, a memorized lookup
table, and a degenerate encoder-decoder. The four-level hierarchy we propose---interchange effectiveness, interchange faithfulness, distributional fidelity,
and structural diagnostics---provides a complete evaluation framework.

Second, nonlinear extensions of causal abstraction require generative structure
and structured sparsity. Unconstrained nonlinear featurizers achieve perfect IIA
vacuously; the Structured pi-SAE's combination of ELBO, label-conditional prior,
and $L_1$ sparsity prevents this failure mode. End-to-end intervention training
closes the gap between training and evaluation objectives, achieving IIA~=~0.97
on GPT-2 hypernymy where DAS gets 0.07. This connects causal abstraction to the
identifiability literature: auxiliary supervision plus generative constraints are
necessary for valid nonlinear causal variable recovery.

Third, continuous distributional metrics (KL, JS, normalized logit difference)
should replace binary IIA as the primary evaluation measure. Binary IIA
saturates too easily and hides important quality differences between methods.

Fourth, the mechanism underlying grokking's creation of causal structure is
more nuanced than previously understood. Nonlinearity is necessary but
attention is not (DMF + ReLU groks 18$\times$ faster). Fourier structure
\emph{follows} generalization rather than driving it, and the gradient
discovers this structure $\sim$10{,}000 epochs before the weights crystallize.
The causal subspace itself is gauge freedom---10 seeds produce orthogonal
subspaces that all support equally valid causal interventions. And spectral
profiling, while informative for single-task models, fails completely under
superposition. These findings suggest that \emph{functional} diagnostics
(equivariance, intervention faithfulness) are more robust than
\emph{structural} diagnostics (spectral profile, stratum classification)
for evaluating causal geometry.

The linearity assumption underlying DAS is not a benign default. It is a testable
hypothesis. And the natural fix---adding nonlinearity---introduces a new failure
mode that is invisible to the standard evaluation metric. This paper provides
the diagnostics to detect both problems, the mechanistic analysis to understand
why they arise, and the structured nonlinear method that avoids them.

\paragraph{Data and code.} All code for model training, DAS fitting, structured
VAE training, geometric diagnostics, and figure generation is available at
\url{[repository URL]}.

\bibliographystyle{tmlr}
\bibliography{references}

\end{document}
